% 第 53 章　从一个电子到元素周期表
\chapter{从一个电子到元素周期表}

第~52~章把氢原子这台“样板间”里里外外看了一遍：一个电子被原子核的库仑引力扣留，量子规则只允许它住进一组离散的态，每个态由四个量子数标记。可是抬头看看周围的世界——你手里的铁勺有二十六个电子，金戒指有七十九个，空气里的氖有十个。如果多电子原子只是把氢原子按比例放大，那么元素越多电子就该越相像，世界应该单调得令人绝望。事实恰恰相反：钠遇水爆炸，它的邻居氖却懒到拒绝和任何东西反应；碳能搭出生命，硅能搭出芯片。这一百来种性格迥异的元素，是从同一套量子规则里长出来的。本章的任务是找到让这一切发生的那两条新规则：一条关于“电子彼此不可分辨”，一条关于“一个量子态只住一个电子”。它们合在一起，把散乱的电子排成了元素周期表。

\step{反常识开场}

现在，把手掌平放在桌面上，用力往下按。桌子纹丝不动。这平常得近乎无聊——直到你想起第~51~章讲过的卢瑟福散射实验：原子内部空旷得离谱，原子核只占原子体积的万亿分之一，其余的空间里只有几片电子云。你的手掌由原子组成，桌面也由原子组成，两团“几乎全是虚空”的东西贴在一起，为什么不是互相穿过去，而是被一堵看不见的墙结结实实挡住？

先别急着说“当然是电荷排斥”。这个说法马上就会在本章被拆穿：手掌和桌面的原子整体都是电中性的，而中性的两团东西靠得再近，长程的库仑力也应该是零，或者因电子云互相拉扯而微微相吸——绝不该是一堵推不倒的墙。可桌子明明硬得能站人。更悬的是：这堵“墙”不只挡住你的手，它还规定了水的不可压缩、金属的坚硬，甚至决定了一颗死掉的恒星能扛住自身引力的碾压。撑起整个物质世界的这股力量，课本上找不到它的名字，因为它根本不是一种新的力。

再看一个同样不起眼、同样解释不了的对比。五金店能买到氖气灯管，通电就亮，但你永远无法让氖和任何东西发生化学反应；而元素表上紧挨着它的钠（氖是十号，钠是十一号，只差一个电子），一小块丢进水里会嘶嘶作响、起火甚至爆炸。一个电子的差别，凭什么造成“绝对懒惰”和“极端暴躁”的天壤之别？周期表里这种脾气每隔几个元素就轮回一次——二、八、八、十八，这些数字像一串密码，至今没有用旧物理给出过任何解释。本章结束时，这串密码会自己开口说话。

\step{先猜一猜}

在往下读之前，把下面三个预测写下来。

第一猜：手掌穿不过桌面，你认为起作用的是（a）原子核互相撞上了，（b）电子与电子之间的静电排斥，（c）某种与电荷无关的新机制？如果你选了（b），请回答：两个电中性的原子靠近时，负电的电子云和正电的原子核同时存在，吸引和排斥为什么不是正好抵消，而偏偏剩下排斥？

第二猜：回忆第~52~章的氢原子。如果给原子核不断增加电荷，电子也一个个多起来，你猜原子的大小会（a）越来越大，（b）越来越小，（c）差不多大、而且呈周期性起伏？再猜电离能（把最外层一个电子拽走所需的能量）会随电子数单调上升，还是忽高忽低？

第三猜，也是最要紧的一猜：周期表的周期长度为什么是 \(2,8,8,18,18,\dots\)？这些数有什么共同点？把它们写成 \(2\times1\)、\(2\times4\)、\(2\times4\)、\(2\times9\)、\(2\times9\)，看出规律了吗？那个乘数 \(2\) 又是从哪里冒出来的？

写好之后夹上书签。本章的每一节都会回来审问其中一猜。

\step{动手实验}

\begin{experiment}[针筒里的两堵墙]
\textbf{材料}：一支干净的塑料注射器（不要针头），水，一点食用盐（可选），燃气灶或蜡烛（可选，需成人陪同）。

\textbf{第一步}：拔掉针头，把注射器抽满一半空气，用手指死死堵住出口，用力推活塞。空气被明显压缩，松手后还会弹回来。记下你大约能把体积压掉多少。

\textbf{第二步}：把注射器抽满水，不留气泡，同样堵住出口，用更大的力气推活塞。这次活塞几乎纹丝不动——不是摩擦力大（轻轻晃动时活塞是顺滑的），而是水本身拒绝被压缩。你可以让整个身体的重量压上去，体积的变化仍然小得看不见。

\textbf{第三步（可选）}：取一小撮食盐撒进燃气灶或蜡烛的火焰，火焰立刻变成明亮的黄色；换一截擦干净的铜丝伸进火焰，颜色变绿。每种元素都在火焰里报出自己的名字。

\textbf{记录与思考}：空气和水都由“几乎全空”的原子或分子组成，为什么一个能压掉一半，另一个连千分之一都压不动？堵住你手指的那堵墙，究竟在针筒的哪个地方？请写下你的解释，并特别回答：如果这堵墙是静电排斥造成的，为什么电中性的水分子之间会有净排斥？
\end{experiment}

这个实验的数字版本我们会在后文细算：水的体积被压小百分之一，需要约 \(2\times10^{7}\)~Pa 的压强，相当于在每平方厘米上站一个两百公斤的人；而你用手推针筒，拼尽全力的压强也比这小两个数量级。至于那撮盐的黄光——它是钠原子最外层那个“暴躁电子”的签名，等本章结束你会明白为什么偏偏是黄色。

\step{旧模型遇到麻烦}

我们手上最好的旧模型有两件：库仑定律，和第~52~章的氢原子量子图像。先用它们认真解释桌子和针筒，看看在哪里翻车。

第一场翻车：中性原子之间不该有这堵墙。手掌和桌面的原子都是电中性的。两个中性原子相距较远时，正电荷和负电荷两两配对，库仑吸引和库仑排斥精确抵消，剩下的只有电子云互相极化产生的微弱吸引（正是这股吸引让水能凝结成液体）。按这个图像，两个原子靠近到电子云开始重叠时，吸引应当先增强；即便继续挤压，经典物理里也没有任何一条定律禁止两团带相反电荷的云互相穿插、滑过。旧模型预言：物质应该像两团烟雾互相穿过，或者像两团软泥越压越扁。针筒里的水却硬得像石头。模型给出的不是定量误差，而是定性的全错。

第二场翻车：多电子原子不该有周期律。把氢原子的图像直接推广：原子核带 \(+Ze\) 的电荷，\(Z\) 个电子全都想待在能量最低的态里。第~52~章说过，最低的态就是最小的那个电子云，半径大约是 \(a_0/Z\)（\(a_0\) 是氢原子半径，约 \(5.3\times10^{-11}\)~m）。于是这个模型预言：九十二号元素铀的原子应该只有氢原子的九十二分之一大，不到一个皮米；而实际测量的铀原子半径约 \(1.5\times10^{-10}\)~m，比氢原子还大三倍左右，误差达到上百倍。电离能的预言同样离谱：如果所有电子都挤在最内层，每种元素的性格应该随着 \(Z\) 单调地变得更“抠门”，电子一个比一个拽不动——可门捷列夫在一八六九年就发现，元素的性质是周期性轮回的：锂、钠、钾像一家人，氟、氯、溴像一家人，每到稀有气体就“全家躺平”。一个单调的模型，生不出一张周期的表。

第三场翻车更深一层：旧模型说不清“两个电子”和“一个电子的两倍”是不是一回事。经典物理里，两颗台球永远可以贴上标签，一号球、二号球，各走各的轨道。可所有对电子的测量——质量、电荷、自旋——都给出完全相同的结果，精确到仪器分辨率的极限，没有人找到过任何一颗“长得不一样”的电子。如果电子真的没有身份证，那么“第一个电子在左边、第二个在右边”和“第二个在左边、第一个在右边”到底是两种状态还是同一种状态？这个听起来像文字游戏的问题，恰恰是周期表的出生证明。

\step{发明新概念}

像历史上的物理学家一样，我们把这个文字游戏当真，发明三样新东西。

第一样：\textbf{全同粒子}。把“电子不可分辨”从测量事实升级为记账规则：交换任意两个电子，得到的不算新状态。注意这条规则比“测不出差别”强得多——它说的是差别在原则上不存在。这立刻改变统计的方式。两只袜子，一黑一白，抽法有四种；两颗可以贴标签的球，交换前后算两种排列；而两个电子的“交换前后”是同一份物理现实，多记一次就是错账。这像什么？像银行账户里的钱：你账上的两百元不分“哪一张”，把两笔一百元互换，账目不发生任何变化。它不像什么？不像两粒外观相同的台球——台球就算涂得一模一样，原则上仍能靠连续盯梢追踪各自的身份，而对电子连“盯梢”这个操作都不存在：你无法在不打扰量子态的情况下给粒子挂上名牌。

第二样：\textbf{两种交换性格}。量子理论允许全同粒子只有两种记账方式。交换两个粒子时，量子态要么原样不变——这类粒子叫\textbf{玻色子}；要么整体翻一个符号——这类粒子叫\textbf{费米子}。一个整体符号听起来无关痛痒（测量只看概率，符号翻不翻，概率一样），但它的后果天翻地覆。而且大自然表现出一条严格的搭配规律：自旋为整数的粒子（光子、氦-4 原子）都是玻色子，自旋为半整数的粒子（电子、质子、中子）都是费米子。这条“自旋—统计搭配”在这里只能作为实验事实接受，它的深层理由要动用到相对论与量子场的联姻，是第~49~章的题材。电子是自旋二分之一的费米子——这半句就是整张周期表的种子。

第三样，本章的主角：\textbf{泡利不相容原理}。把费米子的“交换翻号”规则用到“两个电子占据完全相同的量子态”这个假想情形上：交换这两个电子，等于什么都没换，态应当不变；可费米子规则又要求它翻号。一个东西既要等于自己又要等于自己的相反数，唯一的可能就是它是零——这个假想情形在自然界根本不存在。用原子物理的语言说：\textbf{同一原子中，不可能有两个电子拥有完全相同的四个量子数}。不是“很难”，不是“能量上不划算”，而是量子力学的账目里根本没有这一行。请立刻注意这条原理精确的措辞：它没有说两个电子不能在同一位置（电子云处处可以重叠），它说的是不能占据同一个量子态——第~52~章那套由 \(n\)、\(l\)、\(m_l\)、\(m_s\) 四个量子数编址的“房间”，每间最多住一个电子。

这三样新东西一到位，多电子原子的全部结构立刻有了骨架。电子从能量最低的态开始填，填一个少一间房。主量子数 \(n\) 决定“楼层”，角量子数 \(l\) 决定这一层的“房型”（\(l=0,1,2,3\) 分别记作 s、p、d、f 亚层），磁量子数 \(m_l\) 是同一房型的不同“朝向”，自旋 \(m_s\) 只有两个取值。于是第 \(n\) 层总共 \(2n^2\) 个房间：第一层两间，第二层八间——你在“先猜一猜”里看出的密码 \(2\times1\)、\(2\times4\)、\(2\times9\)，那个乘数 \(2\) 就是电子自旋的两个取向。周期表的前两个周期长度 \(2\) 和 \(8\)，就这样从四个量子数的清点里掉了出来。

但真实的能级次序还要过一道修正：电子多了以后会互相“挡光”。最内层的电子云屏蔽了原子核的一部分正电荷，外层电子感受到的不是 \(+Ze\)，而是一个打了折的\textbf{有效核电荷}。屏蔽的程度依赖轨道形状：s 电子云有一小截“钻”到核附近的概率（这叫钻穿效应），被屏蔽得少，能量压得更低；d、f 电子云几乎完全被内层隔在外面。于是能级不再严格按楼层排队：第四层的 4s 被压到第三层的 3d 之下，填充次序变成 1s、2s、2p、3s、3p、4s、3d、4p……周期长度 \(2,8,8,18\) 里的第二个 \(8\) 和第一个 \(18\)，正是这次插队的结果。

最后是同一亚层内部的排座规则——\textbf{洪特规则}：电子填入同一亚层的几个等价轨道时，先分头各占一个，并且自旋保持平行，万不得已才两两同屋。物理原因有两层，都不神秘：分占不同轨道让电子在空间上错开，彼此之间的库仑排斥能更小；而自旋平行的排法在全同粒子的账本上还会额外降低能量（这叫交换效应，它是量子记账方式的结论，不是一种新发现的力）。洪特规则不是新的基本定律，而是静电排斥加上全同粒子规则在亚层里的推论；它只管一件事——在几种能量几乎相同的填法里，大自然选能量最低的那种。

至此，氦-3 与氦-4 这对著名的“双胞胎反例”也登场了。它们的化学性质几乎完全相同（核外都是两个电子），直觉会说：低温下它们的表现也该差不多。错。氦-4 的原子核含两个质子两个中子，加上两个电子共六个费米子——偶数个费米子拼成的复合粒子整体是玻色子，它在 \(2.17\)~K 以下突然变成超流体，黏滞度消失，能沿杯壁爬出来；而氦-3 少一个中子，五个费米子，整体仍是费米子，要到约 \(2.5\)~mK——温度低近千倍——才经由一种完全不同的配对机制获得超流性。只差一个中子，化学家几乎分不出它们，物理学家却看到两个世界：一个是挤不进同一张椅子的费米子，一个是成千上万个可以叠进同一个量子态的玻色子。你的直觉在这里出错，是因为直觉从不数粒子的个数是奇是偶——而大自然数。

\step{一句话模型}

\begin{oneline}
电子是全同的费米子，每个量子态最多住一个；于是一个原子里的电子从能量最低的态开始逐间占房，占满一层再上楼——被屏蔽和钻穿修正过的能级次序决定了占房的节奏，这个节奏就是元素周期表，而最外层住着几个电子，就是元素全部的化学性格。
\end{oneline}

这个模型像什么？像一座按楼层售票的剧院：一人一座，从最好的位置开始卖，卖完一排再开下一排。它不像什么？有三处必须立刻标明。第一，“座位”不是空间里的格子，而是四个量子数编址的量子态——两个电子可以坐在空间同一点，只要量子态不同。第二，观众没有先到先得的自由：占房顺序由能级次序和洪特规则硬性规定，不是电子“喜欢”这样坐，而是其他任何坐法都不满足量子账本的约束或能量更高。第三，满员的楼层（稀有气体的满壳层）不是“观众不想动”，而是再塞一个电子必须开新楼层，代价高到化学反应出不起——惰性是能量账算出来的，不是脾气。

\step{数学翻译器}

先做纯清点。给定主量子数 \(n\)，角量子数 \(l\) 可以取 \(0,1,\dots,n-1\)；给定 \(l\)，磁量子数 \(m_l\) 可以取 \(-l,\dots,0,\dots,+l\)，共 \(2l+1\) 个；每个轨道再乘自旋的两个取向。于是第 \(n\) 层的房间总数为
\[
N_n=\sum_{l=0}^{n-1}2(2l+1)=2n^2 .
\]
这条式子回答四个问题。它描述什么：一层楼最多容纳多少电子。为什么这样写：\(2l+1\) 是朝向的数目，\(2\) 是自旋的数目，求和是把本层所有房型加在一起——等差数列求和正是初中数学。何时成立：只要每个电子可以用四个量子数近似描述（独立电子近似，下文会审问这个假设）。何时失效：它从不给出错误答案，但对强关联的多电子体系，“每个电子各占一个轨道”这幅图像本身会失真。立刻做特殊检查：\(n=1\) 给出 \(N_1=2\)——第一周期恰有氢、氦两种元素；\(n=2\) 给出 \(N_2=8\)——第二周期八种元素；\(n=3\) 给出 \(18\)，但第三周期只有 \(8\) 种——清点没错，是能级次序插了队：3d 的十个房间被 4s 挤到后面，和第四层一起才开张。公式在两个简单楼层上全对，在第三层上迫使我们把“屏蔽”请回来，这正是好公式的脾气：它对不上的地方，恰是物理藏东西的地方。

再算有效核电荷这笔真账。把最外层电子粗略看成在“核加内层电子云”的库仑场里运动的单电子，第~52~章的氢原子公式就能借用：电离能约为
\[
I\approx 13.6\ \text{eV}\times\frac{Z_{\mathrm{eff}}^2}{n^2},
\qquad Z_{\mathrm{eff}}=Z-\sigma ,
\]
其中 \(\sigma\) 是内层电子屏蔽掉的电荷数。特殊检查先走一遍：\(\sigma=0\)（没有内层电子，就是氢）退回第~52~章的 \(13.6\)~eV，对；\(\sigma=Z-1\)（屏蔽得滴水不漏）给出 \(Z_{\mathrm{eff}}=1\)，外层电子像氢的 excited 电子一样只看到一个单位电荷——这是屏蔽的极限，真实原子落在两者之间。现在代入真实数据。钠是十一号元素，\(Z=11\)，最外层是 3s 电子，\(n=3\)，实验测得第一电离能 \(I=5.14\)~eV。反解：
\[
Z_{\mathrm{eff}}=n\sqrt{\frac{I}{13.6\ \text{eV}}}=3\times\sqrt{\frac{5.14}{13.6}}\approx 1.8,
\qquad \sigma\approx 11-1.8\approx 9.2 .
\]
结论：十个内层电子替 3s 电子挡掉了约 \(9.2\) 个核电荷——屏蔽几乎完全，但不彻底，这正是 3s 电子“钻穿”进内层、多感受到一点正电荷的证据。同一个公式用在锂上（\(Z=3\)，2s 电子，\(I=5.39\)~eV）给出 \(Z_{\mathrm{eff}}\approx1.3\)、\(\sigma\approx1.7\)：两个 1s 电子屏蔽了 \(1.7\) 个电荷，同样近乎完全。两行真实数据，把“内层电子挡光、外层电子看残阳”这幅图像钉死在数量级上。

屏蔽把同一个原子劈成了两个世界。铀（\(Z=92\)）的最内层 1s 电子几乎不被屏蔽，感受到全部 \(92\) 个核电荷，结合能约为 \(13.6\times92^2\approx1.15\times10^5\)~eV——和实验测得的铀 K 壳层结合能（约 \(1.16\times10^5\)~eV）吻合；而它最外层的电子只被几个单位的有效电荷拴着，电离能只有 \(6\)~eV 上下。同一个原子里，最内层与最外层电子的结合能相差近两万倍。化学只和最外面那几个“eV 级”的电子打交道，这就是化学反应的能量尺度，也是为什么一张周期表能装下整个化学。

\begin{figure}[htbp]
\centering
\begin{tikzpicture}[>=Stealth,scale=0.62]
  % 能级横线
  \foreach \y/\lab in {0/{1s}, 1/{2s}, 2/{2p}, 3/{3s}, 4/{3p}, 5/{4s}, 6.6/{3d}, 7.6/{4p}}{
    \draw[thick] (0,\y) -- (3,\y) node[right=2pt]{\lab};
  }
  \node[left] at (0,3.8) {能量};
  \draw[->] (-0.7,-0.3) -- (-0.7,7.9);
  % 填充顺序斜线
  \draw[->,dashed,blue!60!black] (1.5,0) -- (1.5,0.9);
  \draw[->,dashed,blue!60!black] (1.5,1.1) -- (1.5,1.9);
  \draw[->,dashed,blue!60!black] (1.5,2.1) -- (1.5,2.9);
  \draw[->,dashed,blue!60!black] (1.5,3.1) -- (1.5,3.9);
  \draw[->,dashed,blue!60!black] (1.5,4.1) -- (1.5,4.9);
  \draw[->,dashed,blue!60!black] (1.5,5.1) -- (1.5,6.5);
  \draw[->,dashed,blue!60!black] (1.5,6.7) -- (1.5,7.5);
  % 电子占位
  \foreach \y/\num in {0/2, 1/2, 2/6, 3/2}{
    \node at (1.5,\y+0.25) {\scriptsize \num 个电子};
  }
  \node[align=left] at (5.9,2.6) {\scriptsize 到 3p 为止，\\ \scriptsize 恰好十八个电子：\\ \scriptsize 氩的满壳};
\end{tikzpicture}
\caption{多电子原子的“售票图”：横线是能级（未按比例），虚线箭头是填充次序。注意 4s 被钻穿效应压到 3d 之下——周期表第三、四周期的长度差就藏在这个插队里。}
\end{figure}

\begin{mathbridge}
\textbf{交换对称性写成公式。}把两个粒子的联合态记作 \(\psi(\mathbf x_1,\mathbf x_2)\)，全同粒子原理要求交换标签后物理不变：\(\psi(\mathbf x_2,\mathbf x_1)=e^{i\alpha}\psi(\mathbf x_1,\mathbf x_2)\)。交换两次等于回到原状，所以 \(e^{2i\alpha}=1\)，只有 \(e^{i\alpha}=\pm1\) 两种解：取 \(+1\) 的是玻色子，取 \(-1\) 的是费米子。泡利原理随之变成一个代数动作：若两个费米子占据同一单粒子态 \(\phi\)，则 \(\psi=\phi(\mathbf x_1)\phi(\mathbf x_2)\) 在交换下不变号，与费米子规则矛盾，故 \(\psi=0\)。\(N\) 个电子的态必须写成完全反对称的组合，它的行列式形式（斯莱特行列式）有一目了然的性质：两行相同，行列式为零——“一间房住不下两个电子”在数学上就是“行列式有两行相同必为零”。

\textbf{简并压：不相容原理的力学化身。}把 \(N\) 个电子塞进体积 \(V\)，泡利原理迫使它们一个换一个地占据越来越高的动量态，填到的最高能量叫费米能
\[
E_F=\frac{\hbar^2}{2m_e}\left(3\pi^2\frac{N}{V}\right)^{2/3}.
\]
特殊检查：密度趋于零，\(E_F\) 趋于零——没有拥挤就没有排斥，对；密度增大，\(E_F\) 按 \(2/3\) 次幂上升。关键在第三行账：即使温度是绝对零度，这些电子也带着巨大动能，压缩 \(V\) 就必须支付这份能量的上涨——能量对体积的导数就是压强。这就是把桌子顶住、把白矮星顶住的\textbf{简并压}：它不是粒子间某种排斥力，而是“不许同态”这条记账规则换算成的力学账单。
\end{mathbridge}

\begin{challenge}
\textbf{为什么偏偏是“一个态一个电子”？}自旋—统计定理给出深层答案，但它需要相对论性量子场论：在三维空间里，要求微观因果性（类空间隔的测量互不影响）与能量正定同时成立，半整数自旋场必须用反对易规则量子化（费米子），整数自旋场必须用对易规则量子化（玻色子）。换句话说，泡利原理不是量子力学里一条可以单独拆掉的经验附加条款——拔掉它，整套理论的自洽性会跟着塌方。这也是为什么一百年来“寻找违反泡利原理的电子”实验精度一再高（某些实验把违反概率的上限压到 \(10^{-30}\) 以下），却始终一无所获，而物理学家并不意外。

\textbf{交换能：洪特规则的定量出身。}两个电子分占轨道 \(a\)、\(b\) 时，反对称化的波函数使两个电子出现在同一位置的概率严格为零（费米孔），自旋平行的排法因此自动减小库仑排斥能。能量修正里出现一个只依赖轨道重叠的积分（交换积分），自旋平行的状态能量更低。请读准主语：让电子自旋“愿意”平行的不是磁性吸引，而是全同粒子的记账方式改变了静电排斥的账单。交换效应是纯粹的量子现象，经典世界没有任何对应物——它却是铁磁性（第~54~章）和洪特规则共同的根。
\end{challenge}

\step{模型的边界}

老规矩，把本章内容拆成三堆。

\textbf{实验事实}：所有电子的质量、电荷、自旋在测量精度内完全相同；原子的光谱、电离能、化学性质呈周期性，周期长度 \(2,8,8,18,18,32\)；稀有气体的满壳惰性、碱金属的单电子活泼性；氦-4 超流与氦-3 的巨大差别；泡利原理至今没有任何被确认的违反。这一堆与诠释无关。

\textbf{数学模型}：独立电子近似——假定每个电子在“核加其余电子平均场”里各自占据一个四量子数轨道，再把泡利原理当作硬约束逐间填房。这套模型何时成立：解释周期表结构、电离能趋势、化学价，它好得惊人；原子的粗略能量，误差常在可接受范围。何时失效：电子之间的关联强到“平均场”三个字遮不住的时候——某些过渡金属氧化物、重费米子材料里，电子的运动互相纠缠，单电子轨道图像失真，必须动用真正的多体量子力学；精确的定量计算（比如某个电离能的第三位有效数字）也总是要回到多体方法。此外，钻穿、屏蔽这些语言本身就是近似模型的话术，不要反客为主当成基本机制。

\textbf{物理解释}：关于“全同粒子为什么只有两种统计”“自旋为何决定统计”，标准答案是相对论性量子场论的自旋—统计定理（见挑战盒子），这不是悬而未决的诠释之争；但它依赖“三维空间”“洛伦兹不变性”等前提——在二维系统里，全同粒子的交换可以给出任意相位（任意子），这已是实验上观测到的实在。所以“费米子与玻色子二分天下”本身也有适用边界。

还有三条最常见的误用必须点名。第一，泡利原理不是一种新的力。它不像电磁力那样有个平方反比公式；它是一条记账约束，约束的力学后果（简并压）通过“能量随体积变化”显现。说“电子靠泡利力互相排斥”和当年说“向心力是一种新力”犯的是同一种病。第二，壳层模型里的“轨道”仍然不是行星轨道——s 电子云是球对称的概率分布，没有任何东西在绕圈，钻穿说的是概率分布的形状，不是电子“有时候抄近路”。第三，不要把费米子“怕挤”、玻色子“爱扎堆”读成粒子有心理活动：两种性格都是交换对称性的数学结论，光子涌进同一个模式不是被吸引过去，而是全同玻色子的概率账本让“已在场的同伴”提高后来者进入同一态的权重——这个区别在第~55~章讲受激辐射时会变成激光的全部秘密。

\step{回头解释开场现象}

回到那只按在桌上的手掌。手掌下压时，皮肤表层原子的电子云与桌面原子的电子云开始重叠。重叠区域里，两边的电子发现自己想占据的量子态已经被对方占据——泡利原理禁止共住，唯一合法的出路是把一部分电子推向能量更高的空态。这要花钱：压缩一点体积，整个系统的能量就上涨一截，而能量对体积的导数正是压强。于是你感受到一堵墙。这堵墙里没有新的力，出面的仍然是老朋友电磁学（电子云重叠改变静电能量）加一条量子账本（泡利原理把退路堵死）。针筒里的水同理：水分子已经被氢键和范德华吸引压到了平衡间距，再压缩就必须重叠电子云、支付简并能量——体弹模量高达 \(2.2\times10^{9}\)~Pa，你的手劲差着好几个数量级，自然压不动。而空气分子相距十倍于分子尺寸，电子云井水不犯河水，压缩只是把分子间的空旷匀掉一些，花的是热运动的小钱，所以一压就扁。同一个“原子几乎全空”，在两种密度下演出两种力学。

氖与钠的天壤之别也结账了。氖（\(Z=10\)）的电子恰好填满 2p 亚层：满壳，再来的电子只能开第三层，而第三层电子感受到的有效核电荷骤降到 \(1\) 上下——往里塞电子不划算；想把满壳里的电子拽走，又要面对 \(21.6\)~eV 的高电离能。两头堵死，于是氖在化学上躺平。钠（\(Z=11\)）多出的那一个电子被泡利原理孤零零地赶上第三层，内层十个电子把核电荷屏蔽到只剩约 \(1.8\)，\(5.14\)~eV 就能把它摘走——丢一个电子，钠就露出氖的满壳内芯。丢电子极其划算，这就是它遇水暴躁的全部原因：水里的氧乐于接收电子，钠乐于甩卖，一拍即合，反应放出的热量点燃氢气。周期表每一行都以“满壳的懒汉”收尾，以“多一个孤电子的暴脾气”开头——周期律不是元素的日历，而是能级填充的里程碑。

至于盐粒撒进火焰的黄光：火焰的能量把钠那个 3s 孤电子踢到 3p，它落回来时放出能量差约 \(2.1\)~eV 的光子，波长恰是 \(589\)~nm 的黄色。一个被泡利原理赶上三楼的电子，用最便宜的方式宣告自己的存在——这就是开场那个可选实验的完整账目，也是第~55~章光谱学的序曲。

\begin{figure}[htbp]
\centering
\begin{tikzpicture}[>=Stealth,scale=1.0]
  % 原子核
  \draw[thick,fill=red!30] (0,0) circle (0.28);
  \node at (0,0) {\scriptsize \(+11\)};
  % 内层电子云
  \draw[thick,dashed,fill=blue!8] (0,0) circle (0.95);
  \node at (0,-1.25) {\scriptsize 内层 10 个电子：\(-10\)};
  % 外层电子
  \draw[fill=black] (3.2,0) circle (0.06);
  \node[above] at (3.2,0.1) {\scriptsize 3s 电子};
  % 有效电荷箭头
  \draw[->,very thick,orange!80!black] (0.35,0) -- (2.9,0);
  \node[above,orange!80!black] at (1.6,0.05) {\scriptsize 有效吸引约 \(+1.8\)};
\end{tikzpicture}
\caption{屏蔽的账本：钠原子核带 \(+11\)，内层十电子云挡掉约 \(-10\)，最外层的 3s 电子只感受到约 \(+1.8\) 的有效核电荷——它因此成为元素表上最容易被摘走的电子之一。}
\end{figure}

\step{脑洞实验与挑战题}

\begin{thought}[如果电子是玻色子]
把宇宙的账目改一行：让电子变成玻色子，一个量子态想住多少住多少。后果在几分钟内席卷一切：所有原子的全部电子都坍进 1s 轨道，原子半径按 \(1/Z\) 缩小，铀原子比氢原子小九十二倍；没有壳层，没有满壳，没有周期表；原子之间再也没有那堵“不许同态”的墙，分子键合的架构整个改写，固体失去硬度，电子云互相穿插如入无人之境——很可能连“桌子”和“手”的分别都不复存在。再把尺度推到极端：真实的宇宙里，大质量恒星烧尽燃料后，引力把原子压碎，电子被挤到密度约 \(10^{36}\)~m\(^{-3}\) 的绝境，费米能被顶到约 \(10^5\)~eV——化学反应能量尺度的十万倍。正是这份由不相容原理支付的简并压，顶住了整颗恒星的引力，造出地球大小、质量却堪比太阳的白矮星；连这条防线也被压垮（超过约 \(1.4\) 倍太阳质量）后，电子被压进质子变成中子，中子自己的简并压再顶出中子星。你此刻按着桌子的手掌，和白矮星不塌缩的原因，是同一条量子记账规则。
\end{thought}

\begin{puzzles}
\item 不用公式，只用“交换翻号”，论证：两个电子不可能四个量子数全同；再论证：两个电子却可以处在空间同一点。这两句话为什么不矛盾？\emph{提示}：泡利原理约束的是“态”，不是“位置”；位置相同的概率分布完全可以由两个自旋取向相反的电子各占一个轨道提供。把“态”和“位置”分开，是本章最重要的一次概念分拣。
\item 钾（\(Z=19\)）的第十九个电子为什么住进 4s 而不是 3d？如果把它硬塞进 3d，它会“看到”什么样的有效核电荷？\emph{提示}：3d 电子云几乎完全被氩的满壳芯（十八个电子）挡在外面，\(Z_{\mathrm{eff}}\) 接近 \(1\)；4s 电子云有一截钻穿到核附近，感受到明显大于 \(1\) 的吸引，能量反而更低。插队不是规则任性，是屏蔽账算出来的。
\item 氦-3 和氦-4 化学性质几乎相同，低温行为却天差地别。请数清两种原子里费米子的个数，说明复合粒子的统计性格由什么决定，并据此判断：一个碳-12 原子（六质子、六中子、六电子）是费米子还是玻色子？\emph{提示}：逐个合计质子、中子、电子；偶数个费米子拼出的整体是玻色子，奇数则是费米子——交换两个复合粒子等于交换其中每一个费米子，符号只取决于个数的奇偶。
\item 有人说：“泡利不相容原理说明电子之间有一种极强的排斥力，所以桌子才是硬的。”请指出这句话里偷换了什么概念，并用“能量随体积的变化”重述桌子变硬的正确因果链。\emph{提示}：不相容原理是一条状态约束，不含任何力公式；压缩迫使电子升上更高的能态，系统能量 \(E(V)\) 随体积缩小而陡增，压强是 \(-\mathrm dE/\mathrm dV\)——力（压强）是能量账的斜率，不是新力种。这和“向心力不是新力”是同一个教训。
\end{puzzles}
